
\medskip
\phantomsection\label{def:sga3-VIB-6.6}%
\noindent\textbf{Definition 6.6.---}
Let \(G\) be an \(S\)-functor in groups and let \(H\) be a group
sub-\(S\)-functor of \(G\). One says that \(H\) is \emph{invariant}
(respectively, \emph{central}, respectively,
\emph{characteristic}) in \(G\) if
\[
  \underline{\operatorname{Norm}}_G H=G
\]
(respectively, if
\(\underline{\operatorname{Centr}}_G H=G\), respectively, if for
every \(S\)-scheme \(T\) and every automorphism
\[
  a\in\operatorname{Aut}_{T\text{-}\mathrm{gr}}(G_T)
\]
one has \(a(H_T)\subset H_T\)). In other words, for every
\(S\)-scheme \(T\), the subgroup \(H(T)\) of \(G(T)\) is normal
in \(G(T)\) (respectively, central in \(G(T)\), respectively,
invariant under every automorphism of \(G_T\)).

\noindent\emph{N.B.}
If \(H\) is central (respectively, characteristic), then it is
invariant.

\medskip
\phantomsection\label{rem:sga3-VIB-6.7}%
\noindent\textbf{Remark 6.7.---}
Let \(G\) and \(H\) be two \(S\)-groups and let \(u:H\to G\) be
a monomorphism. The group \(H\) is invariant (respectively,
central) in \(G\) if and only if the morphism
\[
  \mu\circ
  \bigl(c\circ\operatorname{pr}_2,\,
        \mu\circ(u\times\operatorname{id}_G)\bigr)
  :H\times_S G\longrightarrow G
\]
(defined by
\[
  (h,g)\longmapsto g^{-1}hg
\]
for \(g\in G(S')\) and \(h\in H(S')\)) factors through \(u\)
(respectively, is equal to \(u\circ\operatorname{pr}_1\)).
The group \(H\) is characteristic in \(G\) if and only if, for
every \(S\)-scheme \(T\) and every \(T\)-group automorphism \(a\)
of \(G_T\), the morphism \(a\circ u(T)\) factors through \(u(T)\).

\medskip
\phantomsection\label{ex:sga3-VIB-6.8}%
\noindent\textbf{Example 6.8.---}
Let \(G\) be an \(S\)-functor in groups. Then
\(\underline{\operatorname{Centr}}G\) is characteristic and
central. If \(G_s\) is representable for every \(s\in S\), then
\(G^0\) is characteristic. This follows from the definitions and
\hyperref[prop:sga3-VIB-3.3]{3.3}.

\medskip
\phantomsection\label{par:sga3-VIB-6.9}%
\noindent\textbf{6.9.---}%
\footnote{\label{note:sga3-VIB-67}Editor's note: this subsection
has been added.}
In [RG71], I, \SGARefLink{cor:sga3-VIA-3.3.3}{3.3.3}, the authors introduce the geometric notion
of a \emph{pure \(S\)-scheme}, which is local on \(S\) for the
étale topology; we refer to \emph{loc.\ cit.} for the precise
definition. We merely record the following points:
\begin{enumerate}
\item[\textup{a)}]\phantomsection\label{point:sga3-VIB-6.9-a}
  (\emph{loc.\ cit.}, Theorem 3.3.5) If \(B\) is a flat,
  finitely presented \(A\)-algebra, then \(B\) is a projective
  \(A\)-module if and only if
  \(\operatorname{Spec}(B)\) is pure over
  \(\operatorname{Spec}(A)\).
\item[\textup{b)}]\phantomsection\label{point:sga3-VIB-6.9-b}
  Consequently, if \(X\to S\) is locally of finite
  presentation, flat, and pure, then \(X\) is essentially free
  over \(S\).
\item[\textup{c)}]\phantomsection\label{point:sga3-VIB-6.9-c}
  (\emph{loc.\ cit.}, 3.3.4\textup{(iii)}) If \(X\to S\) is
  locally of finite type, flat, has geometrically irreducible
  fibers, and has no embedded components, then \(X\) is pure over
  \(S\).
\end{enumerate}
Every group scheme locally of finite type over a field is
Cohen--Macaulay
(cf.\ \hyperref[prop:sga3-VIA-1.1.1]{VI\(_A\), 1.1.1}), and
therefore has no embedded components. We obtain in particular:
\begin{enumerate}
\item[\textup{d)}]\phantomsection\label{point:sga3-VIB-6.9-d}
  Every \(S\)-group scheme \(G\) that is locally of finite
  presentation, flat, and has connected fibers is pure over
  \(S\), and hence essentially free over \(S\).
\end{enumerate}
All the statements of
\hyperref[ex:sga3-VIB-6.2.4-e]{6.2.4\textup{(e)}} may therefore
be used with the results of \textup{(b)} and \textup{(d)} above
taken into account.
